\section{Related work}
\label{sec:related}

\paragraph{Waste and dumping detection.}
Image-level litter detection has TACO \citep{proenca2020taco}. ZeroWaste
\citep{bashkirova2022zerowaste} covers cluttered conveyor-belt scenes rather than
in-the-wild litter. UAVVaste \citep{kraft2021uavvaste} and AerialWaste
\citep{torres2023aerialwaste} cover aerial monitoring, and MARIDA
\citep{kikaki2022marida} and MADOS \citep{kikaki2024mados} cover marine pollution.

Video-based dumping detection now has a public benchmark and an active contest.
The 2026 IWDD contest \citep{bouwmans2026overview} ranked ten teams on the
Mivia-IWDD release by precision, recall and $F_1$, normalised notification delay,
frame rate and peak memory. The larger Mivia-IWDD-500 release
\citep{greco2026iwdd} postdates the contest. Seven entries are published:
3D CNNs with temporal localisation \citep{almorsy2026x3d}, cross-fold transformer
ensembles \citep{kim2026kfold}, multi-head classifiers \citep{scarrica2026garberus},
vision--language temporal reasoning \citep{kairanbay2026vlm}, novelty detection
\citep{zbuce2026novelty}, onset-centric estimation \citep{bhat2026imslab} and a
lightweight two-stage recurrent framework \citep{putzu2026unicass}. These are
strong backbones, and SENTINEL-E does not compete with them. It consumes a frozen
backbone's scores, so better backbones and this work compound. None of these
entries reports a false-alarm rate over a continuous horizon, an average run
length, or a detection delay under continuous monitoring. The contest overview
analyses ROC space, but its false-positive-rate axis is also per clip.

\paragraph{Video anomaly detection benchmarks.}
CUHK Avenue \citep{lu2013avenue} and UCSD Ped2 \citep{mahadevan2010ucsd} are
classical benchmarks, and recent work reports clip- or frame-level ROC-AUC on
them. A reconstruction-based idempotent generative network reports 99.03\% AUC on
Ped2 and 92.40\% on Avenue \citep{dong2025rvadign}. For a survey of deep video
anomaly detection, see \citet{liu2024generalized}. AUC is computed after the fact
on a fixed, labelled test set. It does not say how often a rule alarms on an
unlabelled, continuously monitored stream, which is the quantity our real-data
protocol measures. \citet{rashidi2026benchmarkauc} audit this gap directly.
Eleven off-the-shelf feature sets that average 70.4\% AUC on their training
distribution fall to 49.9\% under cross-dataset shift, which is chance. At
thresholds that a same-dataset evaluation would call operational, false-alarm
rates are around $3.2\times 10^{4}$ per hour. A benchmark AUC does not certify
what happens when the clock keeps running.

Our real-data backbones are classical. Optical flow is estimated by Farneb\"{a}ck's
method \citep{farneback2003}. PatchCore scores each patch by its nearest-neighbour
distance to a memory bank of frozen ImageNet features from a residual network
\citep{roth2022total,he2016deep}.

\paragraph{Online detection with a false-alarm bound.}
The closest prior work is \citet{doshi2021online}. It gives an online detector
for surveillance video with an asymptotic bound on the false-alarm rate. Our
guarantee is non-asymptotic and time-uniform, holding simultaneously at all
monitoring times. Doshi and Yilmaz also treat video anomaly detection as continual
learning \citep{doshi2022rethinking}.

\paragraph{Sequential detection of changes that end.}
Classical quickest-detection theory treats permanent shifts. A sub-literature
treats changes of finite duration. Transient change detection
\citep{guepie2012transient,guepie2017transient} minimises worst-case missed
detection under a false-alarm constraint, using a window-limited CUSUM. The
intermittent-change line \citep{sokolov2023intermittent} treats a change of
unknown duration directly. \citet{fuh2015hmm} use a two-state hidden Markov model
for quickest detection. It is the closest classical relative of the two-state
chain in \cref{eq:episodic}. Epidemic changepoint detection
\citep{fisch2022capa,juodakis2020epidemic} treats a segment that departs from and
returns to a background level. These methods use CUSUM, SPRT or window statistics
with asymptotic or fixed-sample guarantees, and none is time-uniform. What we add
is an anytime-valid, conformal-fed construction of the finite-duration change.
The finite-duration change itself is not new, and \cref{sec:episodic} states that
scope.

Mixing over the trajectories of a two-state chain and collapsing them by a forward
pass belongs to the fixed-share, or switching-expert, family
\citep{herbster1998fixedshare,vanerven2012switch}. \cref{eq:episodic} is a
two-expert case, with one expert held at the neutral bet. That literature does not
give the martingale property under optional stopping, which \cref{thm:episodic}
establishes.

Closer to our architecture are the conformal test martingales of
\citet{vovk2003exchangeability} and their successors. They pair a conformal front
end with a Ville-thresholded martingale for online change and exchangeability
testing, the same two-layer shape as \cref{sec:layer1,sec:episodic}. We know of no
application of e-process or anytime-valid conformal methods to waste, dumping or
pollution surveillance.

\paragraph{Sequential change detection.}
CUSUM \citep{page1954}, Shiryaev's Bayesian statistic \citep{shiryaev1963} and
Shiryaev--Roberts \citep{roberts1966} are the classical procedures. Their
optimality theory is due to \citet{lorden1971} and \citet{moustakides1986}, and
\citet{tartakovsky2014sequential} give a full treatment. Their guarantees are
asymptotic average-run-length statements that presume a correctly specified
pre-change law. Detector scores from real footage are autocorrelated, heavy-tailed
and diurnally non-stationary, so that presumption fails, and both procedures lose
control in our experiments. \citet{shin2024edetectors} give a nonparametric
e-detector framework with non-asymptotic ARL bounds. \citet{saha2026nonpartitioned}
extend it to composite pre-change classes. Their PFA-controlled variant mixes over
start times with weights summing to one, which gives
$\mathbb{P}(\tau<\infty)\le\alpha$. \cref{prop:generalisation} reduces to that
variant at zero episode-end rate.

All of these mix over the changepoint location of a persistent shift. None adds an
episode-end axis with a time-uniform, optional-stopping-safe guarantee. Cost also
separates them. Evaluated exactly over all active starts, a non-partitioned
e-detector costs $O(t)$ per observation \citep{saha2026nonpartitioned}. Under a
geometric prior over start times, the sum telescopes to the $O(1)$ recursion of
\cref{eq:forward} at $\eta=0$. The episodic process keeps that cost while adding
the episode-end axis. That matters when the horizon is a month of video and the
hardware is a camera.

\paragraph{Anytime-valid inference and e-values.}
Testing by betting \citep{shafer2021testing} turns Ville's inequality into a design
principle. \citet{waudby2024betting} develop the betting strategies we adapt.
\citet{wang2022ebh} introduce e-BH, which controls the false discovery rate under
arbitrary dependence and underpins our fleet layer. The closest work in spirit is
E-SHIFT \citep{khan2026eshift}. It detects distribution shift in streaming learning
systems with a sub-exponential e-process on a composite non-conformity score, and
it offers an $\alpha$-mixing extension for correlated streams. We reimplement it
as a baseline and find three differences. First, its log-moment-generating-function
bound is estimated, so by the authors' own account validity is asymptotic in the
calibration size rather than finite-sample. Second, it assumes i.i.d.\ calibration
and deployment scores. That is stronger than the exchangeability our conformal
layer needs, and a diurnal cycle breaks it regardless. Third, its wealth restarts
at frame one after each alert, so evidence arriving hours into a stream must first
repay what was burned earlier. Our restart convention between alarms has the same
effect.

\paragraph{Conformal prediction.}
Split conformal prediction \citep{vovk2005alrw} supplies the distribution-free
p-values we bet on, and \citet{angelopoulos2023gentle} give an introduction.
Conformal p-values also test for outliers \citep{bates2023testing}, and exchangeability
can be relaxed for sequential data \citep{barber2023beyond,xu2023sequential}.
\citet{vovk2012conditional} studies calibration-conditional validity, which our
Beta levels address. \citet{tibshirani2019covariate} handle covariate shift by
likelihood-ratio weighting, which we implement and evaluate. Conformalised
residuals \citep{lei2018distribution} and localised methods \citep{guan2023localized}
pursue conditional coverage by other means. Our Layer 1 adds the sequential
requirement. The p-values must be super-uniform conditionally on the past, not
merely marginally, or the wealth process is not a supermartingale.

\paragraph{Control charts and ROC practice.}
Deployment practice has analogues of these quantities. Control charts are tuned to
an in-control average run length, and their limits are revisited when the process
changes \citep{montgomery2020introduction}. Our time-uniform false-alarm probability
is the horizon-free analogue: it bounds the chance of any alarm at any time, so it
does not depend on how long the camera runs. ROC analysis fixes a per-decision
operating point, a trade-off between true and false positive rates
\citep{fawcett2006introduction}. A budget of false alarms per monitoring period is a
different quantity. It depends on how many decisions are made, which a per-decision
ROC point does not show. \cref{sec:limitations} returns to alarm budgets.
